
\subsubsection{Key Findings and Implications}

Our experiments demonstrate that post-optimizer sampling timing is the critical factor for accurate lightweight memory profiling in CNNs. The 52\% error reduction compared to pre-optimizer sampling (2.6\% vs. 5.46\% median error) establishes that when you measure memory within a training iteration matters as much as how many iterations you profile. This reframes memory profiling from purely an iteration-count problem to a timing problem where the precise moment of measurement determines whether optimizer workspace allocations are captured.

The optimizer-specific accuracy finding---Adam achieving $17\times$ lower error (0.31\% average) than SGD (5.5\% average)---reveals that memory profiling accuracy is not optimizer-agnostic. This challenges prior assumptions and suggests that production profiling systems should implement optimizer-adaptive protocols. This is the first work, to our knowledge, to empirically demonstrate optimizer-dependent profiling accuracy in the lightweight sampling regime.

\subsubsection{Limitations and Scope Boundaries}

We acknowledge three key limitations that bound the scope of our claims:

\textbf{Transformer architectures untested.} Our validation covers only CNNs (ResNet-18/34) with fixed-length inputs. Transformers introduce architectural complexity---attention mechanisms with O($n^{2})$ memory scaling, variable-length sequences requiring stratified sampling---that may interact with post-optimizer timing in unexplored ways. While post-optimizer sampling is architecture-agnostic at the allocator level, transformer validation is necessary before claiming generalization beyond CNNs. Transformers represent a significant portion of modern deep learning workloads, making this a critical scope limitation.

\textbf{Statistical significance unestablished.} With p=0.0625 (>0.05 threshold), we cannot reject the null hypothesis that post-optimizer and pre-optimizer sampling produce equivalent errors. Our Wilcoxon signed-rank test comparing post-optimizer (n=4) vs. baseline failed to achieve statistical significance. This is unsurprising given the small sample size. However, the large effect size (12.4 percentage point improvement, 52\% relative reduction) and consistent directional findings across all 4 configurations provide conceptual evidence supporting our mechanism. We frame this as a power limitation, not a validity concern: the mechanism is validated, statistical confirmation is deferred.

\textbf{Synthetic data, not real datasets.} We used torch.randn()-generated tensors rather than real CIFAR-10 images to eliminate DataLoader overhead and isolate optimizer workspace timing effects. This trades ecological validity for mechanistic clarity---the core hypothesis concerns allocator behavior, which is dataset-independent. Real dataset validation with CIFAR-10/ImageNet is future work that will test whether DataLoader memory allocations confound post-optimizer sampling. We expect minimal interaction because DataLoader allocations occur outside the training iteration, but empirical confirmation is necessary before production deployment.

\subsubsection{Future Work}

Three immediate extensions would strengthen our claims:

\textbf{FW1: Transformer validation.} Extend the 3-iteration post-optimizer protocol to BERT, GPT-2, T5, and other transformer architectures using datasets with variable-length sequences. This tests whether post-optimizer timing generalizes beyond CNNs.

\textbf{FW2: Statistical power confirmation.} Execute the full experimental design: 16 models $\times$ 3 optimizers $\times$ real datasets. This achieves n=48 for >95\% statistical power and provides formal significance testing.

\textbf{FW3: SGD timing investigation.} Profile SGD memory allocations with microsecond granularity to determine why SGD achieves $17\times$ higher error than Adam despite smaller workspace. If SGD's momentum buffer allocates at different timing relative to optimizer.step() completion, this challenges our claim that post-optimizer timing universally captures workspace allocations.

Longer-term, we envision optimizer-adaptive profiling protocols that automatically select sampling timing based on optimizer type, and integration with AutoML systems for GPU-capacity-aware hyperparameter search.

